\section{What the mathematics establishes}
\label{sec:conclusions}

The concept can be expressed as a chain of distinct requirements:
\begin{equation}
 \begin{gathered}
 \text{topology/rule distribution}
 \ \longrightarrow\
 \text{predictive synthetic structure}\\
 \longrightarrow\
 \text{useful backbone displacement}
 \ \longrightarrow\
 \text{language-stage gain}.
 \end{gathered}
\end{equation}
Only the first arrow is fixed by designing a generator.
None of the later arrows is guaranteed by eight-regular connectivity.

\subsection{Results established under stated assumptions}

\begin{enumerate}
 \item A periodic grid NCA is exactly representable by typed message
 passing with matching kernels and relation order
 (Proposition~\ref{prop:convolution}).
 Rewiring is a change of dynamical system, not merely a notation change.

 \item Degree-preserving swaps preserve typed degree and uniform
 endpoint marginals. They do not guarantee connectivity, a larger
 mixing gap, or useful language transfer. An explicit disconnected
 eight-regular graph demonstrates the first obstruction.

 \item A deterministic binary rule has a finite effective truth table.
 For the pilot, complete-trajectory entropy is bounded by 576 bits
 conditional on the graph, despite 1024 raw state bits.
 This does not make gzip ratios universal complexity measures or prove
 that a finite predictor can cheaply exploit the bound.

 \item For full preceding-token windows, direct-parent lags are at most
 33 in the binary pilot. Crop prefixes require a separate qualification.
 Mean structural lag is exactly 18 under every typed-permutation
 rewiring, while grid lag variance is exactly $153/8$.

 \item Favorable backbone displacement satisfies the sufficient
 alignment inequality \eqref{eq:alignment}.
 The fully specified linear surrogate \eqref{eq:exactgraphtarget}
 derives favorable or unfavorable transfer directly from topology and
 target overlap, with both outcomes possible at the same degree.
 A solved quadratic model explains the role of slow error modes;
 a counterexample shows that a lower initial loss need not improve
 subsequent convergence.

 \item Cumulative benefit requires \eqref{eq:breakeven}.
 The explicit $39$-versus-$25$ language-update example produces either
 a total gain or a total loss as synthetic cost changes.
\end{enumerate}

\subsection{The remaining research hypothesis}

A precise, still unproved hypothesis is the following:
\begin{quote}
There exists a declared graph/rule distribution for which, under common
language-stage initialization of the token table, data order and
optimizer, synthetic pre-pre-training reduces language loss or first-hit
cost relative to a matched grid distribution, across independent
rule/topology/data seeds.
\end{quote}
This is an existence and reproducibility claim about a particular
distribution and model family. It is not a theorem about all graph NCAs.
Neither the favorable toy gradients nor the toy quadratic hitting times
are evidence that the project's actual backbone satisfies the required
conditions.

There are several mathematically distinct unresolved quantities:
the rule-conditioned predictive entropy seen through finite windows;
the fraction of loss reducible by markers or repeated patches;
the actual transfer displacement and its language-gradient alignment;
its projection onto slowly decaying language directions;
and the cost needed to obtain it. Better values of one quantity do not
force better values of all the others.

\subsection{How to extend this note}

Future mathematical work can replace individual assumptions explicitly:
use stochastic transition kernels in \eqref{eq:stochasticentropy};
analyze another serialization through \eqref{eq:lag};
derive update-direction bounds for a specified optimizer;
or study a declared graph ensemble rather than assuming that partial
rewiring is a uniform random graph.
Any new theorem should state its quantifiers, regularity assumptions,
and the observable quantity it controls. A new numerical claim should
show the equation and its substituted values, distinguish a constructed
example from a recorded measurement, and retain the failure cases.

No dynamical rollout, optimization simulation, or new experimental run
was used for this note. Its numerical values are evaluations of the
displayed algebra, matrix spectra, counts, and previously recorded
measurements. The document therefore provides a mathematical working
foundation, while leaving the empirical transfer question open.

